\documentclass[11pt]{article}
\usepackage[margin=1in]{geometry}
\usepackage{amsmath,amssymb,amsthm,mathtools}
\usepackage{graphicx}
\usepackage{microtype}
\usepackage[hidelinks]{hyperref}

\newtheorem{theorem}{Theorem}
\newtheorem{lemma}{Lemma}
\newtheorem{remark}{Remark}
\newcommand{\A}{\mathcal A}
\newcommand{\D}{\mathcal D}
\newcommand{\V}{\mathcal V}
\newcommand{\E}{\mathbb E}

\title{A Universal Constant Two for Derivations on In-Tree Tensor Algebras}
\author{Candidate full answer to Conjecture 6.27 of arXiv:2507.22508}
\date{26 August 2026}

\begin{document}
\maketitle

\begin{center}
\fbox{\parbox{0.91\textwidth}{\textbf{Status.}
Candidate full solution, likely valid, pending human review.  For every finite
in-tree $G$ and every derivation $\delta$ on its tensor algebra $\A_G$, there
is $S\in\A_G$ such that $\delta=\delta_S$ and
$\|S\|\leq 2\|\delta\|$.  Thus the conjectured universal constant exists,
with $C=2$.}}
\end{center}

\section{Source conjecture}

Huang and Ma ask for a graph-independent implementation constant for
derivations on tensor algebras of finite in-trees \cite[Conjecture 6.27]{HM}.
Here an in-tree is a finite directed tree with a range root $v_0$ such that
there is a directed path from every vertex to $v_0$.

\begin{figure}[h]
\centering
\includegraphics[width=0.94\textwidth]{figures/open_problem_crop.png}
\caption{Conjecture 6.27 in the source, arXiv PDF page 33.}
\end{figure}

With the source convention, $\delta_T(A)=AT-TA$.  The sign convention is
irrelevant to the norm estimates below.  The source proves that every
derivation on $\A_G$ is inner when $G$ is a finite tree
\cite[Lemma 6.2]{HM}; the open point is a bound independent of the number and
shape of the vertices.

\section{Proof intuition}

Start with any implementer $T$.  Averaging $T$ over the diagonal sign
unitaries separates it into its vertex-diagonal part $D$ and its
positive-length path part $N$.  The averaging identity writes $N$ as an
average of values of the derivation on norm-one diagonal unitaries, so
$\|N\|\leq\|\delta\|$.

The diagonal coefficients are controlled without adding estimates along
edges.  For a whole directed path $p$, the $L_p$ coefficient of
$\delta(L_p)$ is exactly the difference of the diagonal coefficients at the
two endpoints.  Applying this to the path from each vertex directly to the
root bounds the entire diagonal, modulo the root scalar, by $\|\delta\|$.
Adding the two independent pieces gives the constant $2$.

\section{The theorem}

For a path $p$, write $s(p)$ and $r(p)$ for its source and range, and let
$L_p$ be the corresponding left creation operator.  The vertex projections
$(L_v)_{v\in\V(G)}$ are pairwise orthogonal and sum to the identity.

\begin{theorem}\label{thm:main}
Let $G$ be a finite in-tree.  For every derivation
$\delta:\A_G\to\A_G$, there exists $S\in\A_G$ such that
\[
  \delta=\delta_S,\qquad \|S\|\leq 2\|\delta\|.
\]
Consequently Conjecture 6.27 holds with the universal constant $C=2$.
\end{theorem}

\begin{proof}
By the source's Lemma 6.2, choose $T\in\A_G$ with $\delta=\delta_T$.
Let
\[
 \D_G=\operatorname{span}\{L_v:v\in\V(G)\}.
\]
For a sign vector $\varepsilon=(\varepsilon_v)_{v\in\V(G)}\in
\{-1,1\}^{\V(G)}$, put
\[
 U_\varepsilon=\sum_{v\in\V(G)}\varepsilon_vL_v.
\]
This is a unitary element of $\D_G\subseteq\A_G$.  Define the diagonal
expectation
\[
 E_{\D}(X)=2^{-|\V(G)|}\sum_\varepsilon
 U_\varepsilon XU_\varepsilon^*.
\]
If $p$ has positive length, then $r(p)\neq s(p)$ because $G$ is a tree, and
\[
 U_\varepsilon L_pU_\varepsilon^*
 =\varepsilon_{r(p)}\varepsilon_{s(p)}L_p.
\]
The sign average therefore kills every positive-length path term and fixes
every vertex term.  Write
\[
 D=E_{\D}(T)=\sum_{v\in\V(G)}\lambda_vL_v,
 \qquad N=T-D.
\]
Since
\[
 T-U_\varepsilon TU_\varepsilon^*
 =(TU_\varepsilon-U_\varepsilon T)U_\varepsilon^*
 =-\delta(U_\varepsilon)U_\varepsilon^*,
\]
we obtain
\begin{equation}\label{eq:offdiag}
 \|N\|
 \leq 2^{-|\V(G)|}\sum_\varepsilon
       \|\delta(U_\varepsilon)\|
 \leq \|\delta\|.
\end{equation}

Now fix a positive-length path $p$.  The Fourier coefficient at $L_p$ is the
contractive matrix functional
\[
 \phi_p(X)=\langle X\xi_{s(p)},\xi_p\rangle.
\]
The diagonal part gives
\[
 \phi_p(L_pD-DL_p)=\lambda_{s(p)}-\lambda_{r(p)}.
\]
Every term in $L_pN$ or $NL_p$ is either zero or is supported on a path of
length strictly larger than $\ell(p)$.  Thus it has zero $L_p$ coefficient,
and hence
\begin{equation}\label{eq:path}
 |\lambda_{s(p)}-\lambda_{r(p)}|
 =|\phi_p(\delta(L_p))|
 \leq \|\delta(L_p)\|
 \leq \|\delta\|.
\end{equation}

Let $v_0$ be the range root.  For each vertex $v\neq v_0$, the in-tree
property supplies a directed path $p_v$ with $s(p_v)=v$ and
$r(p_v)=v_0$.  Applying \eqref{eq:path} to this entire path gives
\[
 |\lambda_v-\lambda_{v_0}|\leq\|\delta\|.
\]
Since the $L_v$ are orthogonal projections,
\begin{equation}\label{eq:diag}
 \|D-\lambda_{v_0}I\|
 =\max_{v\in\V(G)}|\lambda_v-\lambda_{v_0}|
 \leq\|\delta\|.
\end{equation}
Finally set $S=T-\lambda_{v_0}I$.  Subtracting a scalar does not change the
inner derivation, while \eqref{eq:offdiag} and \eqref{eq:diag} give
\[
 \|S\|\leq\|N\|+\|D-\lambda_{v_0}I\|
 \leq 2\|\delta\|.
\]
\end{proof}

\section{Verification, scope, and novelty boundary}

The companion script \texttt{code/verify\_tree\_averaging.py} constructs the
finite left-regular path matrices for random rooted in-trees and checks the
sign-average identity and the exact path-coefficient formula.  The computation
is diagnostic only; the proof is the finite averaging and coefficient
calculation displayed above.

The theorem proves the existence assertion of Conjecture 6.27.  It does not
assert that $C=2$ is optimal.  The source obtains $C=1$ for out-trees and for
one star-shaped family of in-trees \cite[Propositions 6.25--6.26]{HM}.

Bounded exact-title, arXiv-id, author, exact-conjecture, and keyword searches
through 26 August 2026 found the source and mirrors of it, but no later primary
source recording this proof or a resolution of Conjecture 6.27.  Priority is
not certified and expert review remains necessary.

\section{Human review recommendation}

Check three convention-sensitive points: the formula
$U_\varepsilon L_pU_\varepsilon^*=\varepsilon_{r(p)}
\varepsilon_{s(p)}L_p$; the assertion that products involving $N$ cannot
contribute to the $L_p$ coefficient; and the use of the whole path from a
vertex to the root rather than an edge-by-edge sum.  Once these are verified,
the constant-two estimate follows immediately.

\sloppy
\begin{thebibliography}{9}
\raggedright
\bibitem{HM}
L. Huang and M. Ma,
\emph{Free semigroupoid algebras and the first cohomology groups},
arXiv:2507.22508 (2025).
\end{thebibliography}

\end{document}
